% Built by pdfsmith (template: latex). Edit doc.yaml for the PDF metadata and this file for the content.
% Compiled with tectonic (XeTeX): fontspec, bibtex and package downloads are handled for you.
\documentclass[11pt,a4paper]{article}

\usepackage{fontspec}   % Latin Modern by default. To use your own font, put the file in fonts/ and write
                        %   \setmainfont{Inter.ttf}[Path=fonts/, BoldFont={Inter-Bold.ttf}]
\usepackage{microtype}
\usepackage[margin=25mm]{geometry}
\usepackage{booktabs}
\usepackage{graphicx}
\usepackage{amsmath}
\usepackage{siunitx}
\usepackage[numbers]{natbib}
\usepackage[
  pdftitle={A Short Study of Price Elasticity},
  pdfauthor={Product Team},
  pdfsubject={Sample academic article generated by pdfsmith},
  pdfkeywords={pricing, elasticity, LaTeX},
  pdflang={en-US},
  colorlinks=true, linkcolor=black, citecolor=blue!50!black, urlcolor=blue!50!black,
  bookmarksnumbered=true
]{hyperref}
\usepackage{xcolor}

\title{A Short Study of Price Elasticity}
\author{Product Team}
\date{30 September 2026}

\begin{document}
\maketitle

\begin{abstract}
We estimate how demand responds to small price changes and show how the result
feeds a simple revenue model. The article is a template: it exercises a figure,
a table, equations, units and a bibliography so that you can replace the content.
\end{abstract}

\section{Introduction}
Price elasticity of demand measures the relative change in quantity for a
relative change in price~\citep{varian1992}. Typeset with care, even a short
note reads like a paper~\citep{knuth1984}.

\section{Model}
\label{sec:model}
Let $q(p)$ be demand at price $p$. The point elasticity is
\begin{equation}
  \varepsilon(p) = \frac{\mathrm{d}q}{\mathrm{d}p}\cdot\frac{p}{q(p)}.
  \label{eq:elasticity}
\end{equation}
Revenue $R(p)=p\,q(p)$ is maximised where $\varepsilon(p)=-1$, as follows from
Equation~\ref{eq:elasticity}.

\subsection{Data}
We use a price change of \SI{5}{\percent} and an average basket of \num{12500} items;
Table~\ref{tab:results} lists the estimates.

\begin{table}[htbp]
  \centering
  \caption{Estimated elasticity by category.}
  \label{tab:results}
  \begin{tabular}{lrr}
    \toprule
    Category & Elasticity & Std.\ error \\
    \midrule
    Electronics & $-1.42$ & 0.08 \\
    Grocery     & $-0.61$ & 0.05 \\
    Fashion     & $-1.87$ & 0.12 \\
    \bottomrule
  \end{tabular}
\end{table}

\section{Results}
Figure~\ref{fig:demand} shows a demand curve drawn with plain \LaTeX{} primitives.
Replace it with a real image placed next to this file and loaded with \texttt{\textbackslash includegraphics}.

\begin{figure}[htbp]
  \centering
  \setlength{\unitlength}{1mm}
  \begin{picture}(80,50)
    \put(0,0){\vector(1,0){80}}
    \put(0,0){\vector(0,1){50}}
    \qbezier(5,45)(15,10)(75,4)
    \put(60,-6){Quantity}
    \put(-4,52){Price}
  \end{picture}
  \caption{A downward-sloping demand curve.}
  \label{fig:demand}
\end{figure}

\section{Conclusion}
Elasticity differs strongly across categories (Section~\ref{sec:model}), so a
single price rule is unlikely to fit all of them.

\bibliographystyle{plainnat}
\bibliography{references}

\end{document}
